%=====================================================================
% 13  二维共形场论
% 原文：Section 13. Two-dimensional conformal field theory（arXiv:2409.03117v1, p. 182--205）
%=====================================================================
\section{二维共形场论}

\begin{keytakeaways}
\begin{itemize}
  \item \textbf{物理图像：}二维无质量标量场可分裂为左行波与右行波（全纯与反全纯两部分）；把空间方向紧化为圆周并杀掉零模后，理论成为一列独立谐振子构成的自由量子场论，其算子代数由 Heisenberg 代数与 Virasoro 代数描述。
  \item \textbf{数学工具：}zeta 函数正规化把发散的和 1+2+3+... 正规化为 -1/12，由此定出真空能与中心荷；Dedekind eta 函数使环面配分函数在模变换下具有良好性质；正规序乘积、复合算子与算子乘积展开给出顶点代数结构。
  \item \textbf{关键结论：}自由玻色子的 Fock 表示构成中心荷为 1 的 Virasoro 最高权模；环面配分函数在补上零模后（再除以根号下的虚部因子）是模群不变的模函数，即量子层面仍保有共形对称性。
  \item \textbf{易错点：}Wick 转动后全纯与反全纯的区分、相位因子与动量算子的排序、顶点算子的多值相位（一般不是空间局域的）、eta 函数模变换的权重，以及圆值理论中半径取倒数的对偶与有理性（半径平方为有理数时只有有限个扇区），都是易错之处。
\end{itemize}
\end{keytakeaways}

\bigskip
\noindent\rule{\textwidth}{0.6pt}
\bigskip

\subsection{二维经典无质量标量场} \label{mlsb}
考虑 $\Bbb R^2$ 上带拉格朗日量 $\mathcal L=\frac{1}{2}(d\phi)^2$ 的自由无质量标量玻色子 $\phi$。
此时局部泛函 $\phi(t,x)$ 满足二维波动（即弦）方程
$$
\phi_{tt}-\phi_{xx}=0, 
$$
于是它分裂为两个泛函之和
$$
\phi=\tfrac{1}{\sqrt{2}}\phi_L+\tfrac{1}{\sqrt{2}}\phi_R,
$$
其中
$$
\phi_L(t,x)=\psi_L(x+t),\ \phi_R(t,x)=\psi_R(x-t),
$$
它们因显然的原因被称为\emph{左行波与右行波}。
换言之，我们有
$$
(\partial_t-\partial_x)\phi_L=0,\ (\partial_t+\partial_x)\phi_R=0.
$$
于是得到
$$
\phi_x+\phi_t=\sqrt{2}\psi_L'(x+t),\ \phi_x-\phi_t=\sqrt{2}\psi_R'(x-t). 
$$
因此 $\psi_L',\psi_R'$ 的 Poisson 括号为
$$
\lbrace \psi_L'(x),\psi_L'(y)\rbrace=\delta'(x-y),\ \lbrace \psi_R'(x),\psi_R'(y)\rbrace=-\delta'(x-y),
$$
$$
\lbrace\psi_L'(x),\psi_R'(y)\rbrace=0.
$$

于是作 Wick 转动（把 $t$ 换成 $it$，并使 $\phi$ 变为复值），令 $u:=x+it$，便有
$$
\overline\partial_u \phi_L=0,\ \partial_u \phi_R=0,
$$
亦即 $\phi_L=\psi_L(u)$ 是全纯的，而 $\phi_R=\psi_R(\overline u)$ 是反全纯的。

现在考虑 $x$ 取值于圆周 $\Bbb R/2\pi \Bbb Z$ 的情形，并把 Lebesgue 测度归一化为总体积 $1$。此时若仍要把 $\phi$ 分解为左行波与右行波，就应通过要求 $\int_0^{2\pi}\phi(t,x)dx=0$ 来“杀掉零模”（否则弦方程有解 $\phi(t,x)=t$，它不能写成左行波与右行波周期波之和）。于是可以引入取值于 $\Bbb C^\times$ 的坐标 $z=e^{iu}$，而 $\phi_L,\phi_R$ 成为 $\Bbb C^\times$ 上的全纯、反全纯场，分别记为
$\varphi$、$\varphi^*$。于是有 Laurent 展开
$$
\varphi(z)=\sum_{n\in \Bbb Z}\varphi_n z^{-n},\ \varphi^*(\overline z)=\sum_{n\in \Bbb Z}\varphi^*_n \overline z^{-n},
$$
且 $\varphi_0=\varphi_0^*=0$。
当 $z$ 位于单位圆上时，这些正是
$\phi_L(0,x)$、$\phi_R(0,x)$ 的 Fourier 展开；并且对
$$
a(z)=\sum_{n\in \Bbb Z}a_nz^{-n-1}:=i\partial_z \varphi(z),\ a^*(\overline z)=\sum_{n\in \Bbb Z}a_n^*\overline z^{-n-1}:=-i\overline\partial_z\varphi^*(\overline z),
$$
其中 $a_0=a_0^*=0$，
我们有
$$
za=\partial_u\phi_L=\psi_L'(u),\ \overline z a^*=\overline \partial_u\phi_R=\psi_R'(\overline u).
$$
于是对 $z=e^{iu},w=e^{iv}$ 得到
\begin{equation}\label{ftpb}
\lbrace za(z),wa(w)\rbrace=\delta'(u-v).
\end{equation}
注意
$$
\delta'(u-v)=i\sum_{n\in \Bbb Z}nz^{n}w^{-n}.
$$ 
于是令
$$
\delta(w-z):=\sum_{n\in \Bbb Z}z^nw^{-n-1}
$$
（即分布 $\delta(w-z)$ 在 $(S^1)^2$ 上的 Fourier 展开，其中
$|z|=|w|=1$），就可以把 \eqref{ftpb} 写成
\begin{equation}\label{ftpb1}
\lbrace a(z),a(w)\rbrace=-i\delta'(w-z).
\end{equation} 
按分量写出，它取形式
$$
\lbrace \sum_{m\in \Bbb Z}a_m z^{-m},\sum_{n\in \Bbb Z}a_{-n} w^{n}\rbrace=-i\sum_{n\in \Bbb Z}nz^{-n}w^n.
$$
于是得到
\begin{equation}\label{mlsb1}
\lbrace a_n, a_m\rbrace=-in\delta_{n,-m}.
\end{equation} 
类似地，
\begin{equation}\label{mlsb2}
\lbrace a_n^*,  a_m^*\rbrace=in\delta_{n,-m},
\end{equation}
并且
\begin{equation}\label{mlsb3}
\lbrace a_n, a_m^*\rbrace=0,
\end{equation} 
用生成函数写出即为
\begin{equation}\label{ftpb2}
\lbrace a^*(z),a^*(w)\rbrace=i\delta'(w-z), \lbrace a(z),a^*(w)\rbrace=0.
\end{equation} 
最后，我们用 Fourier（即 Laurent）模 $a_n$ 写出该理论的哈密顿量。
回想一下，按原来的记号它具有形式
$$
H=\frac{1}{2}\int_{\Bbb R/2\pi \Bbb Z} (\phi_t^2+\phi_x^2)dx.
$$
于是有
$$
H=\frac{1}{4}\int_{\Bbb R/2\pi \Bbb Z} ((\overline z a^*-z a)^2+(\overline z a^*+za)^2)dx=\frac{1}{2}\int_{\Bbb R/2\pi \Bbb Z}(\overline z^2a^{*2}+z^2a^2)dx,
$$
亦即
\begin{equation}\label{classham}
H=\sum_{n>0}(a_{-n}a_n+a_{-n}^*a_n^*).
\end{equation}
它满足关系
$$
\lbrace a_m,H\rbrace=-ima_m,\ \lbrace a_m^*,H\rbrace=ima_m^*.
$$

\subsection{$\Bbb R\times \Bbb R/2\pi \Bbb Z$ 上杀掉零模的自由量子无质量标量场}

现在考虑 $\Bbb R\times \Bbb R/2\pi \Bbb Z$ 上自由无质量标量玻色子 $\phi$ 的自由量子场论，取 Minkowski 度规 $dt^2-dx^2$ 并杀掉零模，也就是说，它是第 \ref{mlsb} 小节所述经典场论的量子化。由于这不是向量空间上的理论，它不满足 Wightman 公理。不过，我们可以自然地量子化对易关系 \eqref{mlsb1}、\eqref{mlsb2}、\eqref{mlsb3}（取 $\hbar=1$），把它们替换为
$$
[a_n,a_m]=n\delta_{n,-m},\ [a_n^*,a_m^*]=-n\delta_{n,-m},\ [a_n,a_m^*]=0.
$$
换言之，对 $a(z)=\sum_{n\in \Bbb Z}a_nz^{-n-1}$、$a^*(\overline z)=\sum_{n\in \Bbb Z}a_n^*\overline z^{-n-1}$，我们有
$$
[a(z),a(w)]=\delta'(w-z),\ [a^*(z),a^*(w)]=-\delta'(w-z),\ [a(z),a^*(w)]=0,
$$
它们量子化了方程 \eqref{ftpb1}、\eqref{ftpb2}（这是第 \ref{harosc3} 小节的分析在场论中的推广，只是这里有无穷多个以正整数标记的谐振子）。由此可见，欧氏空间局域性成立。

这表明我们有一个由无穷多个独立谐振子构成的系统。
为用代数语言重述这一点，考虑无穷维 Heisenberg 李代数
$\mathcal A$，其基为 $a_n,n\ne 0$ 与（中心的）$K$，对易关系为
$$
[a_n,a_m]=n\delta_{n,-m}K.
$$
于是我们看到，理论的 Hilbert 空间 $\mathcal H$ 的某个稠密子空间应承载
$\mathcal A$ 的一对相互对易的作用（分别由左行波与右行波给出），
而 $K$ 分别以 $1$ 与 $-1$ 作用（我们把 $\mathcal A$ 的第二个拷贝记为 $\mathcal A^*$）。

下面描述 Hilbert 空间 $\mathcal H$。
注意李代数 $\mathcal A$ 有一个不可约的\emph{Fock 表示}
$\mathcal F$，它由 $\Omega$ 生成，定义关系为
$$
a_n\Omega=0,\ n>0,\quad K\Omega=\Omega.
$$ 
作为向量空间，$\mathcal F$ 就是\emph{Fock 空间}
$$
\mathcal F=\Bbb C[X_1,X_2,...]
$$ 
（其中 $\Omega=1$），算子
$a_{-n}$（$n>0$）以乘以 $X_n$ 的方式作用，而 $a_n$
以 $n\frac{\partial}{\partial X_n}$ 作用。

现在，系统的哈密顿量（为方便起见我们把它重新标度一个因子 $2$）应满足对易关系
$$
[\widehat H,a_n]=-na_n,\ [\widehat H,a_n^*]=na_n^*.
$$
由此可见，若希望 $\widehat H$ 的谱有下界，且若 $\Omega\in \mathcal H$
是 $\widehat H$ 的最低本征向量，则对 $n>0$ 必须有
$$
a_n\Omega=0,\ a_{-n}^*\Omega=0
$$
但在这种情况下，由 $\Omega$ 在 $a_n,a_n^*$ 作用下生成的空间 $\mathcal D$ 必须是李代数 $\mathcal A\oplus \mathcal A^*$ 的不可约表示
$\mathcal F\otimes \mathcal F^*$，其中
$\mathcal F^*:=\Bbb C[X_1^*,X_2^*,...]$，
$a_n^*$ 以乘以 $X_n^*$ 作用，而对 $n>0$ 有 $a_{-n}^*\mapsto n\frac{\partial}{\partial X_n^*}$。

于是空间 $\mathcal D$ 是多项式代数 $\Bbb C[X_j]$
与 $\Bbb C[X_j^*]$ 的张量积。每个代数 $\Bbb C[X_j]$ 上带有正定内积，
使 $X_j^n$ 成为正交基，且 $||X_j^n||^2=j^nn!$；$\Bbb C[X_j^*]$ 类似。
由此在 $\mathcal F,\mathcal F^*$ 与 $\mathcal D$ 上得到正定内积 $\la,\ra$，
关于它 $a_i^\dagger=a_{-i}$、$a_i^{*\dagger}=a_{-i}^*$。
Hilbert 空间 $\mathcal H$ 就是 $\mathcal D$ 关于 $\la,\ra$ 的完备化。


这意味着量子哈密顿量必由公式
$$
\widehat H=\sum_{n>0}(a_{-n}a_n+a_n^*a_{-n}^*)+C
$$
给出，它由经典哈密顿量 \eqref{classham} 量子化而来（注意湮灭算子写在右边，
以确保无穷和有意义）。我们可以把 $\widehat H$ 写成左行部分与右行部分之和：
$$
\widehat H=\widehat H_L+\widehat H_R,
$$
其中
$$
\widehat H_L:=\sum_{n>0}a_{-n}a_n+\tfrac{C}{2},\ \widehat H_L:=\sum_{n>0}a_{n}^*a_{-n}^*+\tfrac{C}{2}.
$$

% [Tutor Note: 第二式左端在原文中仍写作 $\widehat H_L$，应为 $\widehat H_R$（笔误，按规范保留原式）；两半的真空能各为 $C/2$，相加才给出 $C$。另一个需要注意之点是正规序：$\widehat H$ 的无穷和只有在湮灭算子居右（即 $n>0$ 的 $a_n$ 在右）时才有意义，这正是后面 $\zeta$ 正规化要把 $\sum_{n>0} n$ 解释为 $\zeta(-1)$ 的缘由。此外，$a_{-n}^*a_n^*$ 的写法与左行部分的 $a_{-n}a_n$ 并不完全对称（右行部分的湮灭算子应为 $a_n^*$，原文把 $a_n^*$ 写在左边），核对时以 $[\widehat H,a_n^*]=na_n^*$ 为准。]

\subsection{$\zeta$ 函数正规化} 
目前还看不出 $C$ 的正确取值应当是什么。为回答这一问题，回想单个谐振子的哈密顿量是作用在 $\Bbb C[z]$ 上的 $z\partial_z+\frac{1}{2}$。
这提示 $\widehat H$ 的公式应为
$$
\widehat H=\sum_{n>0}(a_{-n}a_n+a_n^*a_{-n}^*+n)=\frac{1}{2}\sum_{n\ne 0} (a_{-n}a_n+a_n^*a_{-n}^*),
$$
这是对 $H$ 作量子化的一个更对称、更自然的公式。然而这个公式没有意义，
因为级数
$$
1+2+3+...
$$ 
是发散的。不过，我们可以用\emph{$\zeta$ 函数正规化}把它正规化。

具体说来，回想 Riemann $\zeta$ 函数由公式
$$
\zeta(s)=\sum_{n=1}^\infty n^{-s}.
$$
定义。众所周知，该函数可亚纯延拓到整个复平面，在 $s=1$ 处有唯一的（单）极点，
并满足函数方程；函数方程说的是，函数 $\pi^{-\frac{s}{2}}\Gamma(\tfrac{s}{2})\zeta(s)$
在变换 $s\mapsto 1-s$ 下对称：
$$
\zeta(1-s)=\pi^{\frac{1}{2}-s}\frac{\Gamma(\tfrac{s}{2})}{\Gamma(\tfrac{1-s}{2})}\zeta(s).
$$
现在自然会定义
$$
C=1+2+3+...:=\zeta(-1).
$$
但 $s=2$ 时的函数方程蕴含
$$
\zeta(-1)=\frac{\pi^{-\frac{3}{2}}}{\Gamma(-\frac{1}{2})}\zeta(2)=-\frac{\pi^{-\frac{3}{2}}}{2\pi^{\frac{1}{2}}}\frac{\pi^2}{6}=-\frac{1}{12}.
$$
所以从这个角度看，自然应取
$$
C:=-\frac{1}{12}
$$

\begin{remark} 回想对整数 ${\rm g}\ge 1$ 有
$$
\zeta(2{\rm g})=(-1)^{{\rm g}+1}2^{2{\rm g}-1}\frac{B_{2{\rm g}}}{(2{\rm g})!}\pi^{2{\rm g}}.
$$
于是 $\zeta$ 的函数方程蕴含
$$
\zeta(1-2{\rm g})=\pi^{\frac{1}{2}-2{\rm g}}\frac{\Gamma({\rm g})}{\Gamma(\frac{1}{2}-{\rm g})}\cdot (-1)^{{\rm g}+1}2^{2{\rm g}-1}\frac{B_{2{\rm g}}}{(2{\rm g})!}\pi^{2{\rm g}}=
-\frac{B_{2{\rm g}}}{2{\rm g}}.
$$
因此 Harer--Zagier 定理可以解释为：亏格 ${\rm g}$ 的曲线模空间的欧拉示性数是 $\zeta(1-2{\rm g})$。
特别地，对 ${\rm g}=1$ 得到 $SL_2(\Bbb Z)$ 的欧拉示性数，即 $-\frac{1}{12}$。
\end{remark} 

% [Tutor Note: $\zeta$ 正规化的实质是“用解析延拓定义发散和”，而不是真的求和：$\zeta(s)$ 的级数只在 ${\rm Re}\,s>1$ 收敛，$C=1+2+3+\cdots:=\zeta(-1)$ 是利用函数方程把 $s\ge 2$ 处的值延拓到 $s=-1$ 的结果（原文用 $s=2$ 的函数方程算 $\zeta(-1)$，并用到 $\zeta(2)=\pi^2/6$ 与 $\Gamma(-\frac{1}{2})=-2\pi^{1/2}$）。物理上 $C$ 就是真空能（Casimir 能量）：每个模式贡献 $\frac{1}{2}\omega_n$，而 $\omega_n=n$，故总零点能为 $\frac{1}{2}\sum n$，按 $\zeta(-1)$ 取值为 $-\frac{1}{24}$，于是 $\widehat H$ 的常数项 $C=-\frac{1}{12}=2\times(-\frac{1}{24})$（左右各一半）。注意“正确”的 $C$ 由下一小节的模不变性唯一确定，$\zeta$ 正规化只是给出它的自然解释；若改用其他的正规化（如截断）会破坏模不变性。]

\subsection{配分函数的模性质} 
值 $-1/12$ 确实被证明是 $C$ 最自然的取值。为看出这一点，让我们回到欧氏号差 $|du|^2=dt^2+dx^2$，并把理论放到复环面 $E=E_\tau=\Bbb C^\times/q^{\Bbb Z}\cong \Bbb R/2\pi T\Bbb Z\times \Bbb R/2\pi \Bbb Z$ 上，其中
$$
T>0,\ \tau=iT,\ q=e^{2\pi i\tau}=e^{-2\pi T}\in (0,1).
$$
此时，正如我们从量子力学所知，应考虑配分函数
$$
Z(\tau):={\rm Tr}(e^{2\pi i\tau \widehat H}).
$$
注意
$$
\widehat H(P\otimes Q)=(\deg P+\deg Q+C)P\otimes Q,
$$
其中 $P\in \mathcal F$、$Q\in \mathcal F^*$，而次数由
$$
\deg(X_n)=\deg(X_n^*)=n.
$$ 
给出。于是有
$$
Z(\tau)=\frac{e^{2\pi i\tau(C+\frac{1}{12})}}{\eta(\tau)^2},
$$
其中
$$
\eta(\tau):=q^{\frac{1}{24}}\prod_{n=1}^{\infty}(1-q^n)
$$
是 Dedekind $\eta$ 函数。
现在回想，$\eta(\tau)$ 是权 $\frac{1}{2}$ 的模形式，即
$$
\eta(-\tfrac{1}{\tau})=\sqrt{-i\tau}\cdot\eta(\tau).
$$
所以配分函数 $Z$ 只在唯一一个 $C$ 值下具有良好的模性质，而这个值恰好是 $-\frac{1}{12}$。
 
让我们解释为什么应当期望 $Z(\tau)$ 具有模性质。

为此注意，理论的拉格朗日量
$$
\mathcal L(\phi)=\frac{1}{4\pi}\int_E (d\phi)^2 =\frac{1}{4\pi}\int_E d\phi\wedge *d\phi
$$ 
是\emph{共形不变的}，因为它纯粹用 Hodge *-算子写出，而后者只依赖于 $E$ 上的共形结构。运动方程亦然，它是 Laplace 方程 $\Delta \phi=0$。
换言之，我们的经典场论是\emph{共形的}。因此可以期望相应的量子理论也是共形的。
这应意味着 $Z(-\tfrac{1}{\tau})=Z(\tau)$，因为复环面 $E_{-\frac{1}{\tau}}$ 与 $E_\tau$ 共形等价。

话虽如此，我们注意到这一模性质只在一个 $\tau$ 的线性因子意义下成立：
事实上我们有
$$
Z(-\tfrac{1}{\tau})=-i\tau Z(\tau). 
$$
原因在于我们杀掉了零模，而实际上本应把它包含进来（毕竟环面 $E_\tau$ 中的空间圈并非 $SL_2(\Bbb Z)$ 不变的，因而“$\phi$ 沿该圈的积分为零”这一条件也不是不变的）。这将在下一小节完成。

% [Tutor Note: 这个“只差一个 $\tau$ 的线性因子”是量子共形反常（Weyl 反常）的典型表现：经典理论共形不变，但量子化必须选取正规化，从而产生依赖于度规的乘性反常。具体到此例，$Z(\tau)=e^{2\pi i\tau(C+\frac{1}{12})}/\eta(\tau)^2$，而 $\eta$ 的模变换 $\eta(-1/\tau)=\sqrt{-i\tau}\,\eta(\tau)$ 给出 $\eta^{-2}$ 的权重为 $-1$；恰好 $C=-\frac{1}{12}$ 使指数因子化为 $1$，故 $Z=\eta^{-2}$ 是权重 $-1$ 的模形式。注意原文的式子 $Z(-\tfrac{1}{\tau})=-i\tau Z(\tau)$ 与 $\eta$ 的变换不相容：由 $Z=\eta^{-2}$ 直接计算应得 $Z(-\tfrac{1}{\tau})=(-i\tau)^{-1}Z(\tau)$（取 $\tau=iT$、$T>0$ 可直接核对，两边都化为 $T^{-1}Z(iT)$）。只有取 $(-i\tau)^{-1}$ 这个因子，下一段的 $\mathcal Z=(-i\tau)^{-\frac12}Z$ 才在 $\tau\mapsto-1/\tau$ 下不变；原文此处疑为笔误（按规范保留原式）。]

\subsection{含零模} 
 零模对应于弦方程的周期解
$\phi(t,x)=\alpha+\mu t$（$\alpha,\mu\in \Bbb R$），当 $\mu$ 非零时它不能分解为
左行波与右行波周期波之和。因此把零模放回去，
相当于把 Hilbert 空间 $\mathcal H$ 换成 $\mathcal H_{\rm full}:=\mathcal H\otimes L^2(\Bbb R)$，其中 $L^2(\Bbb R)$ 是一个量子力学自由无质量粒子的 Hilbert 空间；同时把哈密顿量 $\widehat H$ 换成
$$
\widehat H_{\rm full}:=\widehat H+\widehat \mu^2,
$$
其中 $\widehat\mu$ 是这个量子力学粒子的量子动量算子，在 $L^2(\Bbb R)$ 上以乘以动量 $\mu$ 的方式作用。于是我们可以写
$$
\mathcal H_{\rm full}=\int_{\Bbb R}\mathcal H_\mu d\mu,
$$
其中 $\mathcal H_\mu=\mathcal F_\mu\otimes \mathcal F_\mu^*$，
而 $\mathcal F_\mu=\mathcal F$ 但把 $a_0=0$ 换成 $a_0=\mu$，类似地
$\mathcal F_\mu^*=\mathcal F^*$ 但把 $a_0^*=0$ 换成 $a_0^*=\mu$。
此时我们仍有
$$
\widehat H=\widehat H_L+\widehat H_R,
$$
其中
$$
\widehat H_L=\tfrac{1}{2}a_0^2+\sum_{n>0}a_{-n}a_n-\tfrac{1}{24},\ \widehat H_R=\tfrac{1}{2}a_0^{*2}+\sum_{n>0}a_{n}^*a_{-n}^*-\tfrac{1}{24}.
$$

根据注 \ref{givesen}，当时间在 $\Bbb R/L\Bbb Z$ 上跑动时，这种粒子的配分函数在相差一个标度因子的意义下是 $L^{-\frac{1}{2}}$。因此完整的配分函数应为
$$
\mathcal Z(\tau)=(-i\tau)^{-\frac{1}{2}}Z(\tau). 
$$
于是我们得到真正的模性质：
$$
\mathcal Z(-\tfrac{1}{\tau})=\mathcal Z(\tau).
$$

我们注意到，函数 $\mathcal Z(\tau)$ 可以自然延拓到任意 $\tau\in \Bbb C_+$（不必是纯虚数），它就等于对“非矩形”复环面 $E_\tau$ 的路径积分。
为解释这一点，注意平移群 $\Bbb R/2\pi \Bbb Z$ 自然地作用在我们的时空上，因此应期望它作用在 Hilbert 空间 $\mathcal H$ 上。
该群的无穷小生成元 $D$ 应满足对易关系
$$
[D,a_n]=na_n,\ [D,a_n^*]=na_n^*
$$ 
（它与 $\widehat H$ 的相应关系只差第一式中的符号）。
由于 $D\Omega=0$，可得
$$
D(P\otimes Q)=(\deg P-\deg Q)P\otimes Q,
$$
亦即
$$
D=\widehat H_L-\widehat H_R.
$$ 
设 $s\in \Bbb R$、$\tau:=iT+s$。那么 Feynman--Kac 公式的一个扭化版本表明：给定 $s\in \Bbb R$，有
$$
Z(\tau)={\rm Tr}(e^{-2\pi T\widehat H}e^{2\pi isD})=|q|^{-\frac{1}{12}}{\rm Tr}(q^{\widehat H_L} \overline q^{\widehat H_R}),
$$
其中 $q=e^{-2\pi(T+is)}=e^{2\pi i\tau}$。于是我们仍有
$$
Z(\tau)=\frac{1}{|\eta(\tau)|^{2}}.
$$
因此
$$
\mathcal Z(\tau)=\frac{1}{\sqrt{{\rm Im}\tau} |\eta(\tau)|^2},
$$
它是 $SL(2,\Bbb Z)$ 的（实解析）模函数，即在
$\tau\mapsto \frac{a\tau+b}{c\tau+d}$（$a,b,c,d\in \Bbb Z$，$ad-bc=1$）下不变。
于是这里我们得到了真正的量子共形对称性（因为复环面 $E_\tau$ 的模空间恰为 $\Bbb C_+/SL_2(\Bbb Z)$）。确实，这个函数显然在 $\tau\mapsto \tau+1$ 下对称，而我们已经看到它在 $\tau\mapsto -1/\tau$ 下不变，而这两个变换生成 $SL_2(\Bbb Z)$。

\subsection{柱面与环面上的关联函数} 
我们也可以考虑量子场 $a$ 与 $a^*$ 的关联函数。
它们分别在 $\mathcal F$ 与 $\mathcal F^*$ 中计算，
并且容易用表示论求出。例如，
对 $n>0$ 有 $a_{n}a_{-n}\Omega=n\Omega$，故两点函数为
$$
\la \Omega, a(z)a(w)\Omega\ra=\sum_{n=1}^\infty nz^{-n-1}w^{n-1}=\frac{1}{(z-w)^2}.
$$
更确切地说，该级数只在 $|w|<|z|$ 时收敛，但函数可解析延拓到所有 $z\ne w$。由于我们的理论是自由的，更高阶的关联函数由 Wick 公式给出：

\begin{proposition}\label{corfu} 有
$$
\la \Omega,a(z_1)....a(z_{2k})\Omega\ra=\sum_{\sigma\in \Pi_{2k}}\frac{1}{\prod_{j\in [1,2k]/\sigma}(z_j-z_{\sigma(j)})^2},
$$
而 $2k+1$ 点关联函数为零。
\end{proposition} 

我们注意到，由于 $\mathcal F$ 作为 $\mathcal A$-模由 $\Omega$ 生成，
这些函数把 $a(z)$ 确定为一个局部算子（=量子场）。更一般地，它们确定算子
$a(z_1)...a(z_r)$（当 $z_i\ne z_j$ 时），而由于空间局域性，这些算子关于 $z_1,..,z_r$ 对称。然而，这些算子在
对角线 $z_i=z_j$ 上并非良定义（有极点）。
 
\begin{exercise}
给出命题 \ref{corfu} 的一个直接的代数证明。
\end{exercise}  

\begin{exercise}
用 theta 函数计算量子场 $\widetilde a(z):=za(z)$ 在环面 $E:=\Bbb R/2\pi T\Bbb Z\times \Bbb R/2\pi \Bbb Z$ 上的归一化两点关联函数。

\textbf{提示.} 该关联函数由
$$
\frac{\la \widetilde a(z)\widetilde a(w)\ra_E}{\la \emptyset\ra_E}={\rm Tr}_{\mathcal F}(\widetilde a(z)\widetilde a(w)e^{-2\pi T\widehat H_L}).
$$
给出。
\end{exercise} 

\subsection{无穷小共形对称性：Virasoro 代数} 

我们已经指出，二维自由无质量标量理论在经典意义下是共形不变的，并看到了这一不变性在量子层面得以保存的若干表现（环面上配分函数的模不变性）。然而，要系统地研究共形对称性，
就需要考虑\emph{无穷小共形对称性}，它由“无穷小共形映射”亦即 $\Bbb C^\times$ 上的全纯
向量场给出。

为简单起见，我们考虑多项式向量场 $P(z)\partial_z$，其中 $P$ 是 Laurent 多项式
（这已足够，因为多项式场在适当的拓扑下在所有全纯向量场中稠密）。
这样的向量场构成一个 李代数，称为\emph{Witt 代数}\label{Witt algebra}（物理学文献中也称\emph{无中心
Virasoro 代数}），我们记之为 $W$。$W$ 的一组方便基是 $\lbrace L_n=-z^{n+1}\partial_z, n\in \Bbb Z\rbrace$，它满足对易关系
$$
[L_n,L_m]=(n-m)L_{m+n},\ m,n\in \Bbb Z.
$$

李代数 $W$ 通过自由无质量标量经典场论的对称性作用于该理论，因为
其拉格朗日量是共形不变的。事实上很重要的是，这个作用只是 $\Bbb R$-线性的而非 $\Bbb C$-线性的，这是好事——这意味着我们有复化 $W_{\Bbb C}=W\oplus W^*$ 的作用，
其中 $W^*$ 是反全纯向量场的李代数；换言之，
我们有两个相互对易的 $W$ 作用。

若我们的理论在量子意义下是共形不变的，则李代数 $W\oplus W^*$
应作用在空间 $\mathcal D$ 上，且与 $\mathcal A\oplus \mathcal A^*$ 的作用相容，也就是说，应使
$$
[L_n,a(z)]=z^{n+1}a'(z)+(n+1)z^na(z), 
$$
$$
[L_n^*,a^*(\overline z)]=\overline z^{n+1}a^{*\prime}(\overline z)+(n+1)\overline z^na_*(\overline z),
$$
$$
[L_n^*,a(z)]=[L_n,a^*(\overline z)]=0,
$$
或者按分量写成
$$
[L_n,a_m]=-ma_{m+n},\ [L_m^*,a_n^*]=-ma_{m+n}^*,\ [L_n,a_m^*]=[L_n^*,a_m]=0.
$$
这样的作用存在吗？为弄清这一点，首先注意算子 $L_0, L_0^*$ 与 $a,a^*$ 的对易关系分别与 $\widehat H_L,-\widehat H_R$ 相同。由于 $\mathcal D$ 是不可约 $\mathcal A\oplus \mathcal A^*$-模，这意味着由 Schur 引理必有
$$
L_0=\widehat H_L+C_L,\ L_0^*=-\widehat H_R+C_R
$$
其中 $C_L,C_R$ 为某些常数。这表明 $L_n$ 必须把 $\mathcal F$ 中的分次移动 $n$，$L_n^*$ 与 $\mathcal F^*$ 同理。

现在类比公式
$$
L_0=\sum_{k\ge 1} a_{-k}a_k+{\rm const},
$$ 
对 $n\ne 0$ 定义
\begin{equation}\label{formvir}
L_n:=\frac{1}{2}\sum_{k\in \Bbb Z}a_{-k}a_{k+n}.
\end{equation} 
容易验证，这个算子作用在 $\mathcal F$ 上（从而作用在
$\mathcal D=\mathcal F\otimes \mathcal F^*$ 上）是良定义的，并且满足所需的对易关系
$$
[L_n,a(z)]=-z^{n+1}a'(z)+(n+1)z^na(z),\ [L_n,a^*(z)]=0.
$$
再次利用 $\mathcal D$ 的不可约性与 Schur 引理，我们
看到：若所需 $W$ 的作用确实存在，则 $L_n$ 必须由公式 \eqref{formvir} 给出
（注意这里不能加常数，因为 $L_n$ 必须移动次数）。因此剩下的就是检验
所构造的算子是否满足 $W$ 的对易关系。

先设 $n\ne -m$。此时利用 Jacobi 恒等式，我们看到
算子 $[L_n,L_m]-(n-m)L_{m+n}$ 与 $a,a^*$ 对易，故再次由 Schur 引理
它必为常数；然而由于它移动次数，我们得到所需关系
$$
[L_n,L_m]-(n-m)L_{m+n}=0.
$$
因此剩下要考虑 $n=-m>0$ 的情形。此时同样的论证表明
$$
[L_n,L_{-n}]-2nL_0=C(n), 
$$
其中 $C(n)\in \Bbb C$；如果对所有 $n$ 都有 $C(n)=0$，我们就得到 $W$ 的作用。
下面计算 $C(n)$。为此注意 $[L_n,L_{-n}]$ 作用在 $\Omega$ 上的本征值为 $2nC_L+C(n)$。因此只需算出
这个本征值，即向量 $L_nL_{-n}\Omega$。

用多项式实现，我们有
$$
L_{-n}\Omega=\tfrac{1}{2}\sum_{0<j<n}X_jX_{n-j}.
$$
于是
$$
L_nL_{-n}\Omega=\tfrac{1}{4}\sum_{0<j<n}j(n-j)\tfrac{\partial^2}{\partial X_j\partial X_{n-j}}\sum_{0<j<n}X_jX_{n-j}=
\tfrac{1}{2}\sum_{0<j<n}j(n-j)=\frac{n^3-n}{12}.
$$
所以
$$
C(n)=\frac{n^3-n}{12}.
$$
于是我们看到，我们几乎得到了 $W$ 的作用，但并不完全——无论怎样选取 $C_L$，关于 $n$ 的三次项都会存在（这是量子反常）！取而代之的是 $W$ 的\emph{射影}表示，而它实际上正是 $W$ 的\emph{中心扩张}的表示。这类射影作用在量子力学中其实很常见，因为量子态对应的并非状状态空间中的真正单位向量，而是相差一个相位因子的向量；在量子态（以及量子可观测量）上，对称群的作用是真正的线性作用。典型的例子是\emph{Heisenberg 不确定性关系} $[\widehat p,\widehat x]=-i\hbar$：经典情形下相平面平移的二维群（或李代数）在量子理论中被三维 Heisenberg 群（李代数）取代；又如\emph{自旋现象}：经典转动对称群 $SO(3)$ 在量子理论中被其双覆盖 $SU(2)$ 取代。

这启发下述定义。

\begin{definition} \emph{Virasoro 代数}是 Witt 代数 $W$ 的一维中心扩张，其基为 $L_n,n\in \Bbb Z$ 与
$C$（中心元），对易关系为
$$
[L_n,L_m]=(n-m)L_{m+n}+\frac{n^3-n}{12}\delta_{n,-m}C.
$$
\end{definition} 

% [Tutor Note: $(n^3-n)/12$ 就是 Gelfand--Fuchs 上闭链，它是 Witt 代数一维中心扩张的代表：$\mathrm{H}^2(W;\Bbb C)$ 是一维的，故中心扩张在差一个标度意义下唯一（习题第 2 问即此）。核对时可以取 $n=1$：$C(1)=0$（故 $L_1,L_{-1},L_0$ 张成 $\mathfrak{sl}_2$ 的通常关系）；取 $n=2$：$C(2)=6/12=\frac12$，于是 $[L_2,L_{-2}]-4L_0=\frac12 C$，这与本题中 $[L_n,L_{-n}]-2nL_0=C(n)$ 的约定一致。注意 $L_0$ 的常数项（真空能）在 $\mathcal F$ 上取 $0$（定理中 $L_0=\sum_{k\ge1}a_{-k}a_k$），在 $\mathcal F_\mu$ 上则多出 $\frac12\mu^2$；中心荷 $c=1$ 是自由玻色子的特征数字，后面维拉索罗生成元的能量-动量张量 $T(z)$ 的四阶极点系数 $c/2$ 即由此而来。另注意原文在“若我们的理论在量子意义下是共形不变的”一段中写 $[L_n,a(z)]=z^{n+1}a'(z)+(n+1)z^na(z)$（缺少负号），与后文 $[L_n,a(z)]=-z^{n+1}a'(z)+\cdots$ 及分量式 $[L_n,a_m]=-ma_{m+n}$ 不一致，阅读时应以带负号者为准。]

于是我们得到由 $C$ 张成的一维中心理想 $\Bbb CC\subset {\rm Vir}$，且
${\rm Vir}/\Bbb CC\cong W$。

于是得到

\begin{theorem} 公式
$$
L_0=\sum_{k\ge 1}a_{-k}a_k,\ L_n=\tfrac{1}{2}\sum_{k\in \Bbb Z}a_{-k}a_{k+n},n\ne 0
$$
定义了 ${\rm Vir}$ 在 $\mathcal F$ 上的作用，其中 $C$ 以 $1$ 作用。
\end{theorem} 

容易验证，同一命题对 $a_0=\mu$ 的空间 $\mathcal F_\mu$ 也普遍成立。
唯一的改变是 $L_0$ 多出一项 $\tfrac{1}{2}\mu^2$：
$$
L_0=\tfrac{1}{2}\mu^2+\sum_{k\ge 1}a_{-k}a_k. 
$$

若 $C$ 在 ${\rm Vir}$ 的表示 $\Bbb V$ 上以标量 $c$ 作用（例如在每个不可约表示上都是如此），则称 $\Bbb V$ 的\emph{中心荷}为 $c$。因此 $\mathcal F_\mu$
是中心荷 $c=1$ 的 ${\rm Vir}$ 表示。

类似地，公式
$$
L_0^*=-\tfrac{1}{2}\mu^2-\sum_{k\ge 1}a_{k}^*a_{-k}^*,\ L_n^*=-\tfrac{1}{2}\sum_{k\in \Bbb Z}a_{k}^*a_{-k+n}^*,n\ne 0
$$
定义了 ${\rm Vir}$ 在 $\mathcal F_\mu^*$ 上的作用，其中中心元 $C^*$ 以 $-1$ 作用（即中心荷 $c=-1$）。

于是我们在空间 $\mathcal D=\mathcal F\otimes \mathcal F^*$ 上得到两个相互对易的 $W$ 的射影作用，
它们只在 $W$ 的中心扩张 ${\rm Vir}$ 上才给出通常的线性作用。尽管如此，$W$ 在量子可观测量上的相应伴随作用
是真正的线性作用，因此这个量子场论是\emph{共形的}。

我们注意到，Virasoro 作用保持 $\mathcal F_\mu$ 上的正定 Hermitian 形式，即
$$
L_n^\dagger=L_{-n}.
$$
因此 $\mathcal F_\mu$ 是 ${\rm Vir}$ 的\emph{正能量酉表示}（正能量是指
$L_0$ 可对角化且谱有下界）。

更一般地，我们可以考虑 $\ell$ 个无质量标量场 $\phi_1,...,\phi_\ell$ 的理论。
此时 $\mathcal D=\mathcal F^{\otimes \ell}\otimes \mathcal F^{*\otimes \ell}$，
而 $\mathcal F^{\otimes \ell}$ 是中心荷 $c=\ell$ 的正能量酉 ${\rm Vir}$-模（即
$\ell$ 个 $\mathcal F$ 拷贝的张量积）。

\begin{exercise} 1. 证明 ${\rm Vir}$ 是 $W$ 的非平凡中心扩张（即作为李代数不同构于 $W\oplus \Bbb C$）。

2. 证明 ${\rm Vir}$ 是 $W$ 的泛中心扩张，即 $W$ 的每个由 $\Bbb C$ 给出的非平凡中心扩张
都同构于 ${\rm Vir}$。
\end{exercise} 

\subsection{共形场论中的正规序、复合算子与算子乘积展开} 

下面我们针对贯穿本章的例子——量子的无质量标量玻色子——把第 \ref{compop} 小节关于正规序、复合算子与算子乘积展开的理论在共形场论情形下作一总结。我们已经看到，算子乘积 $a(z)a(w)$ 只在 $w\ne z$ 时有定义，且在 $w=z$ 处有极点，这导致局部算子 $a(z)^2$ 并非良定义。因此，让我们把该算子乘积在 $w=z$ 附近展开为 Laurent 级数，并找出含 $w-z$ 的负幂的奇异部分。
为此考虑差
$$
:a(z)a(w):=a(z)a(w)-\frac{1}{(z-w)^2}.
$$
$a(z)$ 的关联函数公式蕴含
$$
\la \Omega,a(z_1)...a(z_{i-1}) :a(z_i)a(z_{i+1}): a(z_{i+2})...a(z_n)\Omega\ra=
$$
$$
\sum_{\sigma\in \Pi_{2k}: \sigma(i)\ne i+1}\frac{1}{\prod_{j\in \Pi_{2k}/\sigma}(z_j-z_{\sigma(j)})^2}.
$$
注意该函数在 $z_i=z_{i+1}$ 处正则，因此算子
$:a(z)a(w):$ 在 $z=w$ 处正则，亦即对所有 $z,w\in \Bbb C^\times$ 有定义。
这个算子称为 $a(z)$ 与 $a(w)$ 的\emph{正规序乘积}。
特别地，虽然平方 $a(z)^2$ 没有定义，我们却有良定义的正规序平方 $:a(z)^2:$。

用 Laurent 系数表示，
$$
:a(z)a(w):=\sum_{m,n\in \Bbb Z}:a_na_m:z^{-n-1}w^{-m-1},
$$
其中当 $m\ge n$ 时 $:a_na_m:=a_na_m$，
当 $m<n$ 时 $:a_na_m:=a_ma_n$（模的正规序）。
当然，这种排序只在 $m+n=0$ 时才有影响。
特别地，我们看到
$$
\tfrac{1}{2}:a(z)^2:=T(z):=\sum_{n\in \Bbb Z}L_nz^{-n-2},
$$
即 Virasoro 模 $L_n$ 的生成函数。
这个算子称为\emph{（量子）能量-动量张量}。

由此可见，Virasoro 模 $L_n$ 可以看作理论全纯部分中相应无穷小共形对称性的 Noether 荷。
相应的 Noether 流是 $z^{n+1}T(z)$，因为
$$
L_n=\frac{1}{2\pi i}\oint z^{n+1}T(z)dz.
$$
完整理论的 Noether 荷则是 $L_n+ \overline L_n$，相应的流为 $z^{n+1}T(z)+\overline{z^{n+1}T(z)}$。
特别地，哈密顿量 $H$ 在相差一个常数的意义下是
$L_0+\overline{L_0}$，这与公式 \eqref{quanham} 一致。

类似地，我们可以定义多于两个因子的正规序乘积，
$:a(z_1)....a(z_n):$。这可通过对 $n$ 归纳来做。具体说来，我们有
\begin{equation}\label{indu}
:a(z_0)a(z_1)...a(z_n):=a(z_0):a(z_1)...a(z_n):-\sum_{k\in [1,n]}\frac{:\prod_{j\ne k}a(z_j):}{(z_0-z_k)^2}
\end{equation}
容易看出，
算子 $:a(z_1)....a(z_n):$ 没有奇点，且对所有值 $z_1,...,z_n\in \Bbb C^\times$ 良定义。
于是对每个 $r_1,..,r_n$ 都有算子
$$
:a^{(r_1)}(z_1)...a^{(r_n)}(z_n):=\partial_{z_1}^{r_1}...\partial_{z_n}^{r_n}:a(z_1)....a(z_n):
$$
令 $z_1=...=z_n$，就可以对 $a(z)$ 的任意微分多项式 $P$ 定义局部算子 $:P(a)(z):$。
这个局部算子称为\emph{复合算子}，它是经典场论中相应局部泛函 $P(a)(z)$ 的量子化。

\begin{exercise} （态-算子对应）
证明映射 $P\mapsto P(a)(z)\Omega|_{z=0}$
是良定义的，并给出（多项式）局部算子空间
$\mathcal V$ 与 Fock 空间 $\mathcal F$ 之间的同构。
\end{exercise} 

更一般地，反复使用 \eqref{indu}，
我们有
\scriptsize
$$
:a(z_1)...a(z_n):\cdot :a(w_1)...a(w_m):=
\sum_{I\subset [1,n], J\subset [1,m], s: I\cong J}
\frac{:\prod_{i\notin I}a(z_i)\prod_{j\notin J}a(w_j):}{\prod_{i\in I}(z_i-w_{s(i)})^2}.
$$
\normalsize 于是令 $z_i=z,w_j=w$，得到
$$
:a(z)^n::a(w)^m:=\sum_{k=0}^{\min(m,n)}k!\binom{n}{k}\binom{m}{k}\frac{:a(z)^{n-k}a(w)^{m-k}:}{(z-w)^{2k}}.
$$
例如对 $n=m=1$ 就得到熟悉的关系
$$
a(z)a(w)=\frac{1}{(z-w)^2}+:a(z)a(w):=\frac{1}{(z-w)^2}+\text{regular terms}.
$$

更一般地，对 $n=1$ 与任意 $m$ 得到
$$
a(z):a(w)^m:=\frac{m:a^{m-1}(w):}{(z-w)^2}+:a(z)a(w)^m:
$$
$$
=\frac{m:a^{m-1}(w):}{(z-w)^2}+\text{regular terms}.
$$
对 $m=2$ 可写成
$$
a(z)T(w)=\frac{a(w)}{(z-w)^2}+\text{regular terms},
$$
它编码了 $a_i$ 与 $L_j$ 之间的对易关系。

对 $n=2$、$m=2$ 得到
$$
:a(z)^2::a(w)^2:=\frac{2}{(z-w)^4}+\frac{4:a(z)a(w):}{(z-w)^2}+:a(z)^2a(w)^2:=
$$
$$
\frac{2}{(z-w)^4}+\frac{4:a(w)^2:}{(z-w)^2}+
\frac{4:a(w)a'(w)}{z-w}+\text{regular terms}.
$$
这也可以写成
$$
T(z)T(w)=\frac{1}{2(z-w)^4}+\frac{2T(w)}{(z-w)^2}+\frac{T'(w)}{z-w}+\text{regular terms},
$$ 
它编码了 $L_i$ 之间的对易关系。更一般地，在中心荷 $c$ 下
这一关系形如
$$
T(z)T(w)=\frac{c}{2(z-w)^4}+\frac{2T(w)}{(z-w)^2}+\frac{T'(w)}{z-w}+\text{regular terms}.
$$ 

这些是\emph{算子乘积展开}的最简单例子。事实上，我们有下述定理，其
证明留给读者：

\begin{theorem}\label{opethm}
对任意局部算子 $P,Q\in \mathcal V$，存在唯一的有限局部算子序列
$R_{1},...,R_{N}\in \mathcal V$，使得
$$
P(a)(z)Q(a)(w)=\sum_{j=1}^{N}R_j(a)(w)(z-w)^{-j}+\text{\rm regular terms},
$$
其中 $(z-w)^{-j}:=\sum_{k\ge 0} \binom{k+j-1}{j-1}z^{-j-k}w^k$。
\end{theorem} 

注意空间局域性蕴含 $Q(a)(w)P(a)(z)$ 由同一公式给出，只是
$(z-w)^{-j}$ 沿相反方向展开，即
$(z-w)^{-j}:=-\sum_{k<0} \binom{k+j-1}{j-1}z^{-j-k}w^k$。
于是我们有
$$
[P(a)(z),Q(a)(w)]=\sum_{j=1}^{N}\frac{1}{(j-1)!}R_j(a)(w)\delta^{(j-1)}(w-z).
$$
因此定理 \ref{opethm} 给出了关于 $P$ 与 $Q$ 的模之间对易子的信息。
例如，如上所见，
$$
[a(z),a(w)]=\delta'(w-z), 
$$
并且还有
$$
[a(z),T(w)]=a(w)\delta'(w-z),
$$
$$
[T(z),T(w)]=\frac{c}{12}\delta'''(w-z)+2T(w)\delta'(w-z)+T'(w)\delta(w-z),
$$
其中在我们的例子中 $c=1$。

此外，显然可以把定理 \ref{opethm} 的展开唯一地延续到也包含非负次数的项；
也就是说，只需对固定的 $w$ 把正则项按 $z-w$ 展开成 Taylor 级数。例如我们有
渐近展开
$$
a(z)a(w)\sim \frac{1}{(z-w)^2}+\sum_{k=0}^\infty :a^{(k)}(w)a(w):\frac{(z-w)^k}{k!}
$$
因此一般地有
$$
P(a)(z)Q(a)(w)\sim \sum_{j=-\infty}^{N}R_j(a)(w)(z-w)^{-j}.
$$

这个公式称为 $P$ 与 $Q$ 之积的\emph{算子乘积展开}。
算子乘积展开满足一定的公理，
这意味着它在空间 $\mathcal V\cong \mathcal F$ 上定义了称为
\emph{顶点代数}的代数结构（不过我们在这里不讨论它）。

% [Tutor Note: 这里的要点是“同一个 OPE 的两个方向”：$P(a)(z)Q(a)(w)$ 与 $Q(a)(w)P(a)(z)$ 的展开式相同，但 $(z-w)^{-j}$ 按 $|z|>|w|$ 或 $|z|<|w|$ 两种方式展开，两者之差正是 $\delta$ 函数及其导数，于是 OPE 的极点项直接给出模的对易子（这是“径向序”下的标准对应）。核对方法：取 $P=Q=T$，把 $T(z)T(w)$ 的四阶、二阶、一阶极点系数分别读出为 $\frac{c}{2}$、$2T(w)$、$T'(w)$，再与 $[L_m,L_n]$ 的中心扩张项比较，即可验证 $c=1$；注意原文的 $\delta^{(j-1)}(w-z)$ 前的 $\frac{1}{(j-1)!}$ 因子，以及 $\delta'(w-z)=-\delta'(z-w)$ 的符号约定（与 $[a(z),a(w)]=\delta'(w-z)$ 的写法一致）。]

\subsection{顶点算子} 

\emph{顶点算子}由量子化局部泛函 $e^{i\lambda \varphi(z)}$ 得到，其中
$$
\varphi(z)=-i\int a(z)dz=-i(a_0\log z+\sum_{n\ne 0}\tfrac{a_{-n}}{n}z^{n}+a_0^\vee)
$$
而 $a_0^\vee$ 是积分常数（$a_0$ 的对偶变量）。换言之，我们有
$$
e^{i\lambda \varphi(z)}=e^{\lambda \int a(z)dz}=e^{\lambda(a_0\log z+\sum_{n\ne 0}\frac{a_{-n}}{n}z^{n})}e^{\lambda a_0^\vee}.
$$
这个泛函的自然量子化是算子
$$
X(\lambda,z):=:e^{\lambda(a_0\log z+\sum_{n\ne 0}\frac{a_{-n}}{n}z^n)}:e^{\lambda a_0^\vee}=
$$
$$
=e^{\lambda\sum_{n>0}\frac{a_{-n}}{n}z^n}e^{-\lambda\sum_{n>0}\frac{a_{n}}{n}z^{-n}}
z^{\lambda\mu} e^{\lambda\partial_\mu},
$$
由于最后一个因子，它把 $\mathcal F_\mu$ 作用到 $\mathcal F_{\mu+\lambda}$，
作用方式为 $X_0(\lambda,z)z^{\lambda\mu}$，其中
$$
X_0(\lambda,z):=e^{\lambda\sum_{n>0}\frac{a_{-n}}{n}z^n}e^{-\lambda\sum_{n>0}\frac{a_{n}}{n}z^{-n}}.
$$
这里我们在以 $z^\alpha,\alpha\in \Bbb C$ 为基的 $\Bbb C$ 的群代数上工作。

现在注意，若 $[A,B]$ 与 $A,B$ 对易，则由 Campbell--Hausdorff 公式
$$
e^A e^B=e^B e^A e^{[A,B]}, 
$$
并且
$$
[\sum_{n>0}\frac{a_{n}}{n}z^{-n},\sum_{n>0}\frac{a_{-n}}{n}w^{n}]=\sum_{n>0}\frac{z^{-n}w^{n}}{n}=-\log(1-\tfrac{w}{z}).
$$
于是
$$
X_0(\lambda,z)X_0(\nu,w)=(1-\tfrac{w}{z})^{\lambda \nu}:X_0(\lambda,z)X_0(\nu,w):
$$
对 $|w|<|z|$ 成立。于是得到
$$
X(\lambda,z)X(\nu,w)=(z-w)^{\lambda \nu}:X(\lambda,z)X(\nu,w):
$$
对 $|w|<|z|$ 成立，其中正规序把 $\partial_\mu$ 放在 $\mu$ 的右边。
更一般地，我们看到
$$
X(\lambda_1,z_1)...X(\lambda_n,z_n)=\prod_{1\le j<k\le n}(z_j-z_k)^{\lambda_j \lambda_k} :X(\lambda_1,z_1)...X(\lambda_n,z_n):
$$
对 $|z_1|>...>|z_n|$ 成立。
特别地，记 $\mathcal F_\mu$ 的最高权向量为 $\Omega_\mu$，则有
$$
\la \Omega_{\mu+\lambda}, X(\lambda_1,z_1)...X(\lambda_n,z_n)\Omega_\mu \ra=
\prod_{j=1}^n z_j^{\lambda_j\mu}\prod_{1\le j<k\le n}(z_j-z_k)^{\lambda_j \lambda_k}.
$$
对 $|z_1|>...>|z_n|$ 成立。

我们看到，这个关联函数可以解析延拓到对角线 $z_i\ne z_j$ 的补集，但这一延拓一般不是单值的。
换言之，场 $X(\lambda,z)$ 一般并不满足空间局域性。取而代之的是
\begin{equation}\label{commrel}
X(\lambda,z)X(\nu,w)=e^{\pi i\lambda\nu}X(\nu,w)X(\lambda,z),
\end{equation} 
这里的等式是在沿路径解析延拓的意义下理解的：$v:=w/z$ 沿正实数从区域 $|v|<1$ 走到区域 $|v|>1$，并从上方绕开点 $v=1$。
特别地，
$$
X(\lambda,z)X(\lambda,w)=e^{\pi i\lambda^2}X(\lambda,w)X(\lambda,z),
$$
亦即 $X(\lambda,z)$ 具有“统计 $\lambda^2/2$”（其中
统计 $\alpha\in \Bbb R/ \Bbb Z$ 的意思是：
交换两个算子的次序会产生相位因子
$e^{2\pi i\alpha}$；例如统计
$0$ 对应玻色子，统计 $1/2$ 对应
费米子）。

注意若把对易关系 \eqref{commrel} 使用两次，我们得到乘子
$e^{2\pi i\lambda\nu}$；这对应于算子乘积 $X(\lambda,z)X(\nu,w)$ 一般是多值的这一事实。

这是共形场论中出现\emph{辫子}（braiding）的一个例子。具体说来，关系 \eqref{commrel}
称为\emph{辫化空间局域性}（或\emph{辫化对易性}），因为它
可以看作适当辫化张量范畴中的对易性。

另外注意
$$
X'(\lambda,z)=\lambda :a(z)X(\lambda,z):
$$
其中 $X':=\partial_z X$，并且
$$
[a_n,X(\lambda,z)]=\lambda z^nX(\lambda,z).
$$
因此
$$
[L_n,X(\lambda,z)]=z^{n+1}X'(\lambda,z)+\frac{\lambda^2}{2}
(n+1)z^nX(\lambda,z),
$$
这意味着 $X(\lambda,z)$ 的自旋为 $\lambda^2/2$。
于是 $X(\lambda,z)$ 满足自旋-统计性质，它推广了通常的自旋-统计定理：自旋模 $\Bbb Z$ 等于统计。

如注 \ref{anyons} 所述，这类量子场称为“任意子”（anyons，因为它们可以具有任意的自旋与统计），并且只能存在于二维。对任意子而言最一般的
自旋-统计性质说：若 $X,Y$ 是自旋为 $s_X,s_Y\ge 0$ 的任意子，则
$$
X(z)Y(w)=e^{2\pi i \sqrt{s_Xs_Y}}Y(w)X(z).
$$
特别地，我们看到：若 $\lambda^2\in \Bbb Z$ 为奇数，则 $X(n\lambda,z)$
彼此之间对奇数 $n$ 表现为费米子、对偶数 $n$ 表现为玻色子（也就是说，若 $n_1n_2$ 为偶数，相应的算子 $X(n_1\lambda,z)$ 与
$X(n_2\lambda,z)$ 对易；若 $n_1n_2$ 为奇数，则反对易），而当 $\lambda^2$ 为偶数时，它们全都表现为玻色子（即算子彼此对易）。

\subsection{圆值理论} 

现在考虑 $\Bbb C^\times$ 上取值于圆周 $\Bbb R/2\pi r\Bbb Z$ 的无质量标量理论。这个理论
与线值理论相同，只差零模，而零模带来如下圆值弦方程解：
$$
\phi(t,x)=\alpha+\mu t+Nrx,
$$
其中 $\alpha\in \Bbb R/2\pi r\Bbb Z$、$\mu\in \Bbb R$，而
$N$ 是整数（绕数）。这类解构成的空间
是按 $N$ 的取值标记的若干柱面 $T^*S^1$ 的不交并。
于是在量子理论中得到 Hilbert 空间
$$
\mathcal H_r^\circ=\bigoplus_{N,\ell\in \Bbb Z}\mathcal H_r^\circ(N,\ell), 
$$
其中
$\mathcal H_r^\circ(N,\ell)$ 是
$\mathcal F_{\frac{1}{\sqrt{2}}(\ell r^{-1}+Nr)}\otimes \mathcal F_{\frac{1}{\sqrt{2}}(\ell r^{-1}-Nr)}^*$ 的完备化。因此我们得到环面 $E_\tau$ 上配分函数的如下公式：
$$
\mathcal Z_r^\circ(\tau)=|\eta(\tau)|^{-2}\vartheta_r(\tau,\overline \tau),
$$
其中
$$
\vartheta_r(\tau,\overline \tau):=\sum_{\ell,N\in \Bbb Z}e^{\frac{1}{2}\pi i\tau (\ell r^{-1}+Nr)^2-\frac{1}{2}\pi i\overline \tau (\ell r^{-1}-Nr)^2}=
$$
$$
\sum_{\ell,N\in \Bbb Z}e^{-\pi (\ell^2 r^{-2}+N^2r^2){\rm Im}\tau+2\pi i\ell N {\rm Re}\tau}.
$$
这表明一个有趣的对偶性 $\mathcal Z_r^\circ(\tau)=\mathcal Z_{r^{-1}}^\circ(\tau)$；事实上，
我们看到参数为 $r$ 的整个理论与参数为 $r^{-1}$ 的理论等价。
这个对偶性称为\emph{T 对偶}，它在弦论中起重要作用。

我们还注意到，$\vartheta_r$ 是权 $1$ 的实模形式：
$$
\vartheta_r(-\tfrac{1}{\tau},-\tfrac{1}{\overline \tau})=|\tau| \vartheta_r(\tau,\overline \tau),
$$
这导致函数 $\mathcal Z_r(\tau)$ 的模不变性，正如共形场论所预期的那样。
为看出这一点，只需注意到指数中有一个 $\Bbb Z^2$ 上的二次型，其矩阵为
$$
Q(\tau)=\begin{pmatrix} r^2{\rm Im\tau} & -i{\rm Re}\tau\\  -i{\rm Re \tau} & r^{-2}{\rm Im \tau}\end{pmatrix} 
$$
于是
$$
Q(\tau)^{-1}=|\tau|^{-2}\begin{pmatrix} r^{-2}{\rm Im\tau} & i{\rm Re}\tau\\  i{\rm Re \tau} & r^{2}{\rm Im \tau}\end{pmatrix}=
\begin{pmatrix} r^{-2}{\rm Im\tau'} & -i{\rm Re}\tau'\\  -i{\rm Re \tau'} & r^{2}{\rm Im \tau'}\end{pmatrix}=SQ(\tau')S,
$$
其中 $\tau':=-\frac{1}{\tau}$，$S=\begin{pmatrix} 0& 1\\ 1& 0\end{pmatrix}$。因此结果可由 Poisson 求和公式推出。

我们看到，若 $r^2=\frac{p}{q}\in \Bbb Q$（既约分数），
则这个共形场论具有一个称为\emph{有理性}的特殊性质：Hilbert 空间 $\mathcal H_r^\circ$ 是有限个“扇区”之和
$\oplus_{i=1}^n \mathcal V_i\otimes \mathcal V_i^*$ 的完备化，其中
左行场作用在 $\mathcal V_i$ 上，
右行场作用在 $\mathcal V_i^*$ 上，从而
$\vartheta_r(\tau,\overline \tau)$ 以及 $\mathcal Z_r^\circ(\tau)$
都是全纯函数与反全纯函数之积的有限和（事实上容易看出 $n=2pq$）。
例如，
真空向量 $\Omega$ 包含在
张量积 $\mathcal V(pq)\otimes \mathcal V(pq)^*$ 中，
其中对 $s\in \Bbb Z_{>0}$ 我们定义
$\mathcal V(s):=\oplus_{m\in \Bbb Z}\mathcal F_{m\sqrt{2s}}$。
空间 $\mathcal V(s)$ 是一个顶点代数，
称为偶格 $\sqrt{2s}\Bbb Z$ 所附的\emph{格顶点代数}\label{lattice vertex algebra}。
这个代数由顶点算子
$X(m\sqrt{2s},z)$ 生成（如我们所知，它们满足空间局域性的玻色版本）。

\begin{example} 考虑 $r=1$ 的情形。此时我们有两个扇区：
真空扇区 $\mathcal V(2)\otimes \mathcal V(2)^*$
与另一个扇区 $\mathcal W\otimes \mathcal W^*$，
其中 $\mathcal W=\oplus_{n\in 2\Bbb Z+1}\mathcal F_{\frac{n}{\sqrt{2}}}$。
对应 $\mathcal F_{\frac{n}{\sqrt{2}}}$（$n$ 为奇数）的粒子
是统计为 $\frac{1}{4}$ 的任意子，故它们
满足形如 $X(z)Y(w)=iY(w)X(z)$ 的辫化对易关系。

不难证明，顶点算子
$X(\sqrt{2},z)$ 与 $X(-\sqrt{2},z)$ 的 Fourier 模在 $\mathcal V(2)=\Bbb L_0$ 上、以及在 $\mathcal W=\Bbb L_1$ 上生成李代数 $\mathfrak{sl}_2[z,z^{-1}]$ 的射影作用，而它们恰好是仿射 Kac--Moody 代数 $\widehat{\mathfrak{sl}}_2=\mathfrak{sl}_2[t,t^{-1}]\oplus \Bbb CK$ 的不可约可积表示（即 $\mathfrak{sl}_2[t,t^{-1}]$ 在水平 $k=1$ 处的泛中心扩张，其中 $K$ 以 $1$ 作用）：具体地，
$X(\sqrt{2},z)$、$X(-\sqrt{2},z)$、$\sqrt{2} a(z)$ 分别给出流 $e(z)$、$f(z)$ 与 $h(z)$，其中对 $b\in {\mathfrak{sl}}_2$ 有
$$
b(z):=\sum_n (b\otimes t^n)z^{-n-1}.
$$
这就是所谓的水平 1 不可约可积模的\emph{Frenkel--Kac 顶点算子构造}\label{Frenkel-Kac vertex operator construction}（对任意有限维单边（simply-laced）单李代数都有定义），此处是最简单的特例 $\g=\mathfrak{sl}_2$。
因此，$r=1$ 时自由玻色子的圆值理论就是所谓的\emph{Wess--Zumino--Witten 模型}在李代数 $\mathfrak{sl}_2$ 与水平 $1$ 这一最简单例子中的情形。
\end{example} 

\begin{example} 设 $r=\sqrt{2}$。此时我们有四个扇区：
$\mathcal V_j\otimes \mathcal V_{-j}^*$，$j=0,1,2,3$，
其中 $\mathcal V_j=\oplus_{n\in 4\Bbb Z+j}\mathcal F_{\frac{n}{2}}$。
特别地，$\mathcal V_0=\mathcal V(4)$，而
$\mathcal V_2=\mathcal F_1$ 中的粒子是玻色-费米对应中出现的费米子。
\end{example} 

\subsection{自由无质量费米子} 

与自由无质量玻色子类似，
可以描述自由无质量费米子 $\xi(z)$ 的理论。
如第 \ref{ferlag} 小节所解释的，在二维情形考虑
取值于 ${\rm Spin}(2)=U(1)$ 的恒等（tautological）表示、动能为 $(\xi,\bold D \xi)$ 的手征旋量是有意义的。
于是我们有一个量子场
$$
\xi(z)=\sum_{n\in \Bbb Z+\frac{1}{2}}\xi_n z^{-n-\frac{1}{2}}
$$ 
及其共轭量子场 $\xi_*(\overline z)$。
$\xi(z)$ 的模满足关系
$$
[\xi(z),\xi(w)]_+=\delta(z-w), 
$$
其中 $[,]_+$ 是超对易子。
这给出 Clifford 代数关系
$$
\xi_n\xi_m+\xi_m\xi_n=\delta_{m,-n}
$$
对 $m,n\in \Bbb Z$ 成立。该代数有唯一的不可约正能量表示
$\Lambda=\wedge (\xi_{-1/2},\xi_{-3/2},...)$，其中 $\xi_j$ 对 $j<0$ 以乘法作用、
对 $j>0$ 以微分作用。$\Lambda$ 上存在不变的正定 Hermitian 内积，
使 $\xi_j$（$j>0$）的 Clifford 单项式构成正交基（不变性指
$\xi_j^\dagger=\xi_{-j}$）。于是理论的 Hilbert 空间是 $\mathcal D:=\Lambda\otimes \Lambda^*$ 的完备化，其中 $\Lambda^*$ 是对应反全纯场的 $\Lambda$ 的对偶。

哈密顿量 $H$ 应满足对易关系
$$
[H,\xi_n]=-\xi_n,\ [H,\xi_n^*]=\xi_n^*,
$$ 
于是我们有
$$
H=H_L+H_R,
$$
其中
$$
H_L=\sum_{n>0}n\xi_{-n}\xi_n
$$
而 $H_R$ 类似。
Virasoro 代数
由
$$
L_m=\frac{1}{2}\sum_{n\in \Bbb Z+\frac{1}{2}}n:\xi_n\xi_{-n+m}:,
$$
定义，亦即 $H_L=L_0$。

\begin{exercise} 证明这些算子 $L_n$ 满足中心荷 $c=\frac{1}{2}$ 的 Virasoro 对易关系。
\end{exercise}

\bigskip
\noindent\rule{\textwidth}{0.6pt}
% [Tutor Note: 本章小结——二维自由场论的“共形”结构可以概括为三层。(1) 运动学：无质量标量分解为全纯与反全纯两支，Laurent 模 $a_n$ 满足 Heisenberg 关系的两个拷贝；Wick 转动把 $\Bbb R\times S^1$ 变成复环面 $E_\tau=\Bbb C^\times/q^{\Bbb Z}$，配分函数成为 $\tau$ 的函数。(2) 量子化与正规序：用 $\zeta$ 函数正规化定出真空能 $C=-\frac{1}{12}$，从而 $\eta(\tau)^{-2}$ 的模性质只在 $c=1$ 时成立；补上零模的 $(-i\tau)^{-\frac{1}{2}}$ 因子后得到真正模不变的 $\mathcal Z(\tau)=(\sqrt{{\rm Im}\tau}\,|\eta(\tau)|^2)^{-1}$，这体现了 Weyl 反常（乘性反常）的抵消。(3) 对称性：无穷小共形变换给出 Witt 代数，其唯一的中心扩张就是 Virasoro 代数（Gelfand--Fuchs 上闭链 $(n^3-n)/12$），$c=1$ 的自由玻色子与 $c=\frac12$ 的自由费米子都是其酉最高权模；正规序与 OPE 给出 $T(z)T(w)$ 的四阶极点系数 $c/2$，把这一代数结构抽象出来就是顶点代数。圆值理论进一步给出按绕数与动量格分层的扇区分解：$T$ 对偶 $r\to1/r$ 是格对偶的体现，$r^2\in\Bbb Q$ 时有理性保证只有有限个扇区，$r=1$ 与 $r=\sqrt{2}$ 分别对应 $\widehat{\mathfrak{sl}}_2$ 水平 $1$ 的 WZW 模型与玻色-费米对应；原文在费米子一节把 Clifford 关系中的指标写成 $m,n\in\Bbb Z$（应为 $\Bbb Z+\frac{1}{2}$），阅读时需注意。]

